\subsection{Mixture-of-Experts Routing}

We evaluate whether adaptive expert weighting can reduce the oracle gap observed on the phone-captured dataset.
Figure~\ref{fig:moe_improve} compares three configurations: (i) uniform averaging of the heterogeneous ensemble, (ii) the strongest single backbone (ViT-B/16), and (iii) the learned mixture-of-experts (MoE) router.

Uniform averaging treats all experts equally, regardless of input characteristics.
Under real-world shift, this strategy underperforms the strongest single model, indicating that naive aggregation does not fully exploit architectural diversity.
Replacing uniform averaging with a learned routing mechanism improves accuracy by approximately 2.4 percentage points relative to simple averaging.
This gain is achieved without increasing model capacity or retraining the backbone networks.
The improvement therefore reflects more effective utilization of existing expert diversity.

Despite this improvement, routed performance remains substantially below oracle accuracy.
A significant portion of the ensemble's theoretical upper bound remains unrealized.

\paragraph{Routing Behavior Analysis.}
To understand how the router achieves its improvement, we analyze the distribution of learned gating weights.
Rather than producing a balanced mixture of experts, the router assigns dominant probability mass to a single expert (see Figure~\ref{fig:moe_weights}).
Approximately 85\% of the routing weight is allocated to the ViT backbone, with the remaining CNN experts contributing minimally.
This behavior indicates that the routing mechanism does not learn fine-grained, input-dependent specialization across experts.
As shown in Figure~\ref{fig:moe_missed}, 60.1\% of samples are classified correctly by the chosen expert.
Missed oracle cases account for 12.4\% of samples, representing the only region where improved routing could increase performance without modifying the ensemble.
In contrast, 27.5\% of samples correspond to hard failures, where no expert predicts correctly.

\paragraph{Class-level Routing Errors.}
Furthermore, we examine missed oracle cases in which the router selects the ViT backbone despite another expert predicting correctly.
Figure~\ref{fig:vit_fails} shows that these missed opportunities are not uniformly distributed across classes.
They are concentrated in specific categories, including \textit{mouse}, \textit{notebook}, and \textit{teapot}.
This class-dependent pattern indicates that routing errors are structured rather than randomly distributed across the label space.

\paragraph{ViT Failure Analysis.}
To analyze routing blind spots, we compare image characteristics between samples where the ViT backbone predicts correctly and samples where it fails.
Figure~\ref{fig:vit_bubble} shows that ViT failure cases exhibit systematically higher blur levels.
On average, failure images are approximately 32\% blurrier than successful predictions.
Failure cases also exhibit lower contrast.
Importantly, many of these blurry samples are correctly classified by CNN-based experts.
For example, DenseNet-121 and EfficientNet-B0 jointly recover a substantial fraction of ViT failures on the \textit{mouse} class.
These observations indicate that routing errors are not caused by a lack of alternative experts.
Instead, viable CNN predictions exist in regimes where the ViT backbone degrades under reduced image quality.
